\begin{figure*}[t]
    \centering
      \includegraphics[width=0.9\textwidth]{statistic.pdf}
      \vspace{2mm}
    \caption{Comprehensive Statistic of LRV-Instruction. In (d), \textcolor{blue}{BLUE} means existent object manipulation. \textcolor{pink}{PINK} means nonexistent object manipulation. \textcolor{green}{GREEN} means knowledge manipulation.}
    \label{fig:LRV_stat}
    \vspace{-0.2in}
\end{figure*}

CHAIR \citep{rohrbach2018object} was introduced to evaluate object hallucination in image captioning tasks. However, it 
usually demands complex human-crafted rules. Alternatively, \citep{li2023evaluating, fu2023mme} formulate the evaluation of hallucination as a binary classification task that prompts LMM to output "Yes" or "No". However, it is hard to evaluate the LMM output in an open-ended manner. In addition, both methods highly depend on human-annotated groundtruth answers.


To this end, we introduce \textit{GPT4-Assisted Visual Instruction Evaluation (GAVIE)} as a more flexible and robust approach to evaluate object-level hallucination. The general prompt we use is shown in %Fig. \ref{fig:prompt_eval} in 
the Appendix. GPT4 takes the dense captions with bounding box coordinates as the image content and compares human instructions and model response. Then we ask GPT4 to work as a smart teacher and score (0-10) students' answers based on two criteria. \textit{(1) Accuracy: whether the response is accurate concerning the image content}. \textit{(2) Relevancy: whether the response directly follows the instruction}. We use GPT4-32k-0314 in the experiments. Fig. \ref{fig:model} successfully points out that \textit{"dog, ball"} is inconsistent with the image, and the response from the MiniGPT4 did not address the instruction. Unlike previous evaluation methods \citep{li2023evaluating, rohrbach2018object}, \textit{GAVIE} does not require human-annotated groundtruth answers and can freely adapt to diverse instruction formats. As for the knowledge level hallucination or images that are not from the Visual Genome dataset, we use the groundtruth answers as a reference and compare them with predictions (Fig. \ref{fig:chart_evaluation} in the appendix).


\vspace{-0.1in}
\begin{table}[t]
\begin{subtable}[t]{0.3\textwidth}
\small
\begin{tabular}{lcc}
\toprule[1.5pt]
Backbone & Perception & Cognition\\
\midrule
Original MiniGPT4 & 616.41 & 232.71 \\
\textbf{Finetuned MiniGPT4} & \textbf{895.96} & \textbf{296.43} \\
Original mPLUG-Owl & 967.34 & 276.07 \\
\textbf{Finetuned mPLUG-Owl} & \textbf{1298.78} & \textbf{328.21} \\
\bottomrule[1.5pt]
\end{tabular}
\label{tab:stat1}
\end{subtable}
\hspace{1.2in}
\begin{subtable}[t]{0.3\textwidth}
\setlength\tabcolsep{4.7pt}
\centering
\small
\begin{tabular}{lcc}
\toprule[1.5pt]
Backbone & Acc(Pos) & Acc(Neg)\\
\midrule
Original MiniGPT4 & 0.53 & 0.54 \\
\textbf{Finetuned MiniGPT4} & \textbf{0.58} & \textbf{0.68} \\
Original mPLUG-Owl & 0.62 & 0.55 \\
\textbf{Finetuned mPLUG-Owl} & \textbf{0.69} & \textbf{0.78} \\
\bottomrule[1.5pt]
\end{tabular}
\label{tab:stat2}
\end{subtable}
\vspace{2mm}
\caption{Zero-shot multimodal evaluation on MME \citep{fu2023mme} of MiniGPT4-7B, mPLUG-Owl-7B between original models and LRV-Instruction-finetuned models. The left chart shows perception and cognition scores. The right chart shows the accuracy on the positive set and the negative set.
}
\label{tab:MME}
\vspace{-0.05in}
\end{table}

\vspace{-0.1in}
\begin{table}[t]
\begin{subtable}[t]{0.2\textwidth}
\small
\begin{tabular}{lcc}
\toprule[1.5pt]
Model & Acc & F1\\
\midrule
mPLUG-Owl-7B & 0.52 & 0.68 \\
LLaVA-13B & 0.50 & 0.66 \\
% \midrule
MiniGPT4-13B & 0.73 & 0.71\\
% \midrule
InstructBLIP-13B & 0.86 & 0.87 \\
% \midrule
\textbf{Ours-7B} & \textbf{0.86} & \textbf{0.88} \\
\bottomrule[1.5pt]
\end{tabular}
 \caption{Random Set.}
\label{tab:stat1}
\end{subtable}
\hspace{0.8in}
\begin{subtable}[t]{0.2\textwidth}
\setlength\tabcolsep{4.7pt}
\centering
\small
\begin{tabular}{lcc}
\toprule[1.5pt]
Model & Acc & F1\\
\midrule
mPLUG-Owl-7B & 0.57 & 0.66 \\
LLaVA-13B & 0.50 & 0.66 \\
% \midrule
MiniGPT4-13B & 0.67 & 0.67\\
% \midrule
InstructBLIP-13B & 0.71 & 0.76 \\
% \midrule
\textbf{Ours-7B} & \textbf{0.73} & \textbf{0.79} \\
\bottomrule[1.5pt]
\end{tabular}
\caption{
Popular Set.}
\label{tab:stat2}
\end{subtable}
\hspace{0.7in}
\begin{subtable}[t]{0.2\textwidth}
\setlength\tabcolsep{4.7pt}
\centering
\small
\begin{tabular}{lcc}
\toprule[1.5pt]
Model & Acc & F1\\
\midrule
mPLUG-Owl-7B & 0.60 & 0.64 \\
LLaVA-13B & 0.50 & 0.66 \\
% \midrule
MiniGPT4-13B & 0.62 & 0.63\\
% \midrule
InstructBLIP-13B & 0.63 & 0.72 \\
% \midrule
\textbf{Ours-7B} & \textbf{0.65} & \textbf{0.73} \\
\bottomrule[1.5pt]
\end{tabular}
\caption{Adversarial Set.}
\label{tab:stat2}
\end{subtable}
\vspace{2mm}
\caption{Zero-shot object hallucination evaluation on POPE \citep{li2023evaluating}. Objects not existing in the image are sampled with three different strategies. Random: random sampling, Popular: top-k most frequent objects in MS-COCO, Adversial: objects are first ranked based on co-occurring frequencies, then top-k frequent ones are sampled. \textit{Ours-7B} means \textit{Finetuned mPLUG-Owl-7B}.
}
\label{tab:pope}
\vspace{-0.2in}
\end{table}



\begin{table}[t]
\setlength\tabcolsep{3pt}
\centering
\small
\begin{tabular}{lcccccc}
\toprule[1.5pt]
& \textbf{Ours} & MiniGPT4 & LLaVA & InstructBLIP & MMGPT &mPLUG-Owl\\
\midrule
\textsc{GAVIE-Accuracy (0-10)} & \textbf{6.58} & 4.14 & 4.36 & 5.93 & 0.91 & 4.84\\
\textsc{GAVIE-Relevancy (0-10)} & \textbf{8.46} & 5.81 & 6.11 & 7.34 & 1.79& 6.35\\
\midrule
\textit{Human Expert1 (1-4)}& \textbf{3.48} & 2.61 & 2.87 & 3.00 & 1.90& 2.90\\
\textit{Human Expert2 (1-4)} &\textbf{3.58}& 2.23& 2.07& 2.48& 1.05& 2.27\\
\textit{Human Expert3 (1-4)} & \textbf{3.33} & 2.58 & 2.89 & 2.94 & 1.38& 2.91\\
\bottomrule[1.5pt]
\end{tabular}
% \vspace{0.05in}
\vspace{2mm}
\caption{
Comparison results on our evaluation set evaluated by \textit{GAVIE}. \textit{Ours} means \textit{Finetuned mPLUG-Owl-7B}. All the LMMs are 7B versions to make a fair comparison. 
}
\label{tab:comparison_evaluation}
\end{table}







